\section{数据类型版图：LLM、世界模型、VLA 与机器人}

访谈中多次提到大语言模型、世界模型、VLA、机器人大脑和物理世界 AI。这些概念容易混在一起。可以把它们放到数据类型版图中理解：文本/代码数据适合 LLM，多模态数据提供感知，机器人轨迹提供动作，仿真数据提供可控场景，失败-成功对提供高信息密度反馈。

\begin{figure}[H]
\centering
\includegraphics[width=0.96\textwidth]{data-taxonomy.png}
\caption{数据类型版图：文本、代码、多模态、机器人轨迹、仿真、失败-成功对。}
\end{figure}

\begin{knowledgebox}{读图：不同模型需要不同数据}
LLM 依赖文本和代码；世界模型需要物理预测相关数据；VLA 需要视觉、语言和动作；机器人策略需要端侧轨迹和反馈。数据荒不是绝对没有数据，而是缺少和目标能力匹配的数据。
\end{knowledgebox}

\subsection{术语消化：模型与数据}

\noindent\begin{tabularx}{\textwidth}{p{0.22\textwidth}p{0.32\textwidth}X}
\toprule
术语 & 解决的问题 & 数据需求 \\
\midrule
LLM & 语言理解、推理、代码和数字世界任务 & 文本、代码、多模态语料、反馈数据。 \\
World Model & 预测物理世界状态和未来变化 & 视频、仿真、物理交互、3D/空间数据。 \\
VLA & Vision-Language-Action，连接感知、语言和动作 & 视觉、语言指令、动作轨迹、环境反馈。 \\
Physical AI & 能在物理世界行动的 AI 系统 & 本体数据、仿真、真实部署、评测。 \\
Behavior Benchmark & 面向具身任务的长程评测集 & 仿真环境、任务定义、成功判定。 \\
\bottomrule
\end{tabularx}

\subsection{共生关系}

大模型团队、世界模型团队和 VLA 团队并不是孤立的。VLA 可能使用基础 LLM，世界模型可以作为云端大脑，VLA 可以把行动反馈给世界模型。若评测体系逐渐一致，世界模型和 VLA 的边界也可能变得更接近。

\begin{warningbox}{不要把世界模型和 VLA 混为一谈}
世界模型更关注理解和预测物理世界；VLA 更关注在物理世界采取行动。二者会共生，但目标不同、数据不同、评估方式也不同。
\end{warningbox}

\subsection{本章小结}

数据版图决定模型版图。LLM、世界模型、VLA 和机器人策略需要不同类型的数据，也通过评测和反馈逐渐连接。

